\section{Conclusion}
\label{sec:conclusion}

We have introduced StabilitySection, an analysis framework for measuring the
cross-census stability of GeoSection-optimal redistricting maps, and the
Census Stability Score (CSS), a composite metric that quantifies this stability
along three dimensions: partisan seat outcomes, natural ratio selection, and
proportionality gap magnitude.

The central claim of this paper is that census stability is a qualitatively
distinct form of algorithmic robustness from within-year seed stability.
B.7 establishes that GeoSection finds the same answer regardless of random
initialisation.
T.8 defines how to test whether it finds the same answer when the decennial
census snapshot changes.
The combination is a strong geographic-determinism record when both tests pass:
the observed partition is less likely to be an artefact of random chance,
particular census-vintage data, or one run of the solver.

\paragraph{Key findings.}
From the 2020 GeoSection baseline and the available 2000--2010 proxy table, we
establish the following working picture:

\begin{itemize}
\item \textbf{Several states are candidates for high CSS} (full three-census):
  the Midwest row (IA, KS, NE, IN, MO), the traditional South (AL, AR, LA, KY,
  TN), and New England (ME, WV, RI).
  These states' geographic structures appear stable in the proxy evidence; their
  natural ratios are anchored by physical features---agricultural uniformity,
  mountain topography, fixed city boundaries---that do not change with
  decennial population counts.

\item \textbf{A second group is likely to show moderate stability}:
  the industrial Midwest (IL, PA, MI, OH), the growing South (VA, NC, SC),
  and the Pacific Coast (OR, WA).
  These states often have stable first-level ratios but may see slight seat-count
  shifts as suburban growth continues.

\item \textbf{Fast-growth states are the main volatility candidates}:
  the Sun Belt growth corridor (TX, AZ, FL, GA, NV, CO).
  These states have undergone the most dramatic population redistribution
  between 2000 and 2020; their metropolitan geographies have reorganised in
  ways that shift the natural ratio.
  This would not be an algorithm failure---it would be a signal that the geographic
  structure of these states has genuinely changed.

\item \textbf{The Lorenz drift proxy is useful for triage}:
  states where $p^*$ drifts by $>0.03$ between available estimates are
  flagged for priority re-evaluation at each redistricting cycle.
  This provides a pre-census early warning candidate that does not require
  running the full GeoSection sweep.
\end{itemize}

\paragraph{CSS as a litigation tool.}
A court presented with a GeoSection-optimal map can be told: ``This is the
natural geographic partition of this state.
We know it is natural because three independent censuses---across twenty years
of population change, two political realignments, and a national demographic
transformation---all converge to the same answer.
Any deviation from this map requires an explanation that is not geography.''

For states with CSS $\geq 0.90$, this could become a strong procedural-fairness
record in post-\emph{Rucho} redistricting litigation.
It does not require proof of partisan intent (constitutionally difficult).
It does not require a partisan-symmetry measurement (empirically contested).
It requires the CSS score, the three-census GeoSection run, archived plan
packages, and the jurisdiction-specific legal record.

\paragraph{Future work.}
\begin{itemize}
\item \textbf{Complete three-census CSS.}
  The next evidence lift is the package-complete three-census CSS table:
  archived summaries, assignments, crosswalks, and outcome diagnostics for all
  states where the comparison is meaningful.
  This is the next session's primary task.

\item \textbf{Jaccard district-assignment stability.}
  Once the cross-year tract alignment is implemented, the full district
  Jaccard computation will provide RQ3 (district-assignment stability)
  and RQ5 (Lorenz curve drift analysis) from the plan.

\item \textbf{Population-only perturbation decomposition.}
  Fixing the 2020 graph and substituting 2000 population weights, for the
  10-state decomposition subset, will isolate the contribution of tract
  boundary redesign to instability.

\item \textbf{Predictive model for CSS.}
  Fitting a CSS regression on population growth, boundary growth, seat-count
  change, and $p^*$ will enable pre-census CSS prediction for 2030
  redistricting once the package-complete three-census dataset exists.

\item \textbf{Temporal sensitivity using ACS 5-year estimates.}
  Testing whether $\Delta p^*$ between consecutive ACS releases predicts
  ratio shifts at the next decennial would allow annual CSS monitoring
  rather than decennial-only evaluation.
\end{itemize}

\paragraph{Code and data availability.}
The construction side uses the \texttt{bisect} Rust binary with
\texttt{bisect state --year \{2000,2010,2020\} --partition-mode geosection}
or the multi-year \texttt{BISECT run} workflow.  The CSS aggregation layer
described here is not yet a dedicated production CLI command; the remaining
Phase~2 package should archive the exact run summaries, assignment files,
crosswalks, and scripts used to compute CSS.
